\section{Results}

\subsection{Equilibrium Distribution of Two Symmetric Species}

We first consider the symmetric case of the two-species system, in which $D = 1$, $r = 1$, and $a_{1,2} = a_{2,1} = a$. Under these assumptions, the system can be written in the simplified nondimensional form
\begin{equation*}
    \begin{cases}
        \partial_t N_{1} = \nabla^2 N_1 + N_1(1 - N_1 - a N_2),\\
        \partial_t N_{2} = \nabla^2 N_2 + N_2(1 - N_2 - a N_1).
    \end{cases}
\end{equation*}

We focus on the regime $a > 1$, for which the local dynamics are bistable and the two species cannot coexist in a homogeneous local community. Nevertheless, spatial coexistence may still emerge at the global scale because the two species are dynamically symmetric and neither has an intrinsic advantage. In this sense, the system resembles the Allen--Cahn model, as both exhibit two stable states connected by a transition front. We therefore seek an approximate analytical description of the spatial interface separating the two species. Throughout this section, we assume that species $1$ occupies the left side of the domain, whereas species $2$ dominates the right side.

To simplify the analysis, we introduce the assumption that the total population density remains approximately constant across space. Biologically, this assumption reflects the fact that both species share the same carrying capacity, and competitive interactions suppress growth when the total population approaches that capacity. Define
\begin{equation*}
    u = N_1 + N_2,
    \qquad
    v = N_1 - N_2.
\end{equation*}
The variables $u$ and $v$ represent the total population density and the population contrast between the two species, respectively. Substituting these definitions into the governing equations yields
\begin{equation*}
    \begin{cases}
        \partial_t u = \nabla_x^2 u + u - (N_1^2 + N_2^2) - 2a N_1 N_2,\\
        \partial_t v = \nabla_x^2 v + v - (N_1^2 - N_2^2).
    \end{cases}
\end{equation*}

Using the identities
\begin{equation*}
    u^2 + v^2 = 2(N_1^2 + N_2^2), \qquad
    u^2 - v^2 = 4N_1N_2, \qquad
    uv = N_1^2 - N_2^2,
\end{equation*}
which follow directly from the linear transformation between $(N_1,N_2)$ and $(u,v)$, the system can be rewritten as
\begin{equation*}
    \begin{cases}
        \partial_t u
        = \nabla_x^2 u
        + u
        - \frac{1+a}{2}u^2
        - \frac{1-a}{2}v^2,\\
        \partial_t v
        = \nabla_x^2 v
        + v
        - uv\\
        \hspace{1.35cm}
        = \nabla_x^2 v + v(1-u).
    \end{cases}
\end{equation*}

We now analyze the dynamics of the total density variable $u$. Under the assumption that the total population density varies weakly in space and time, we approximate
\begin{equation*}
    \partial_t u \approx 0,
    \qquad
    \nabla_x^2 u \approx 0.
\end{equation*}
The equation for $u$ then reduces to the algebraic relation
\begin{equation*}
    \frac{1+a}{2}u^2
    - u
    + \frac{1-a}{2}v^2
    = 0.
\end{equation*}
The algebraic relation for $u$ can be solved explicitly in terms of $v$,
\begin{equation*}
    u = \frac{1 \pm \sqrt{1 + (a^2-1)v^2}}{1+a}.
\end{equation*}
Since $u = N_1 + N_2 \geq 0$, only the positive branch is physically meaningful, yielding
\begin{equation*}
    u = \frac{1 + \sqrt{1 + (a^2-1)v^2}}{1+a}.
\end{equation*}
Substituting this expression into the equation for $v$ gives
\begin{equation*}
    \partial_t v
    =
    \nabla^2 v
    +
    v
    \left(
        1
        -
        \frac{1+\sqrt{1+(a^2-1)v^2}}{1+a}
    \right).
\end{equation*}

The effective reaction term
\begin{equation*}
    f(v)
    =
    v
    \left(
        1
        -
        \frac{1+\sqrt{1+(a^2-1)v^2}}{1+a}
    \right)
\end{equation*}
has three equilibria located at $v=0,\pm1$, and is antisymmetric with respect to the origin. The qualitative shape of this function resembles the Allen--Cahn reaction term
\begin{equation*}
    g(v)=rv(1-v^2).
\end{equation*}
To derive an effective approximation, we expand $f(v)/v$ around $v=0$:
\begin{align*}
    \frac{f(v)}{v}
    &=
    1
    -
    \frac{1+\sqrt{1+(a^2-1)v^2}}{1+a}\\
    &=
    \frac{a}{1+a}
    -
    \frac{1}{1+a}
    \sqrt{1+(a^2-1)v^2}\\
    &=
    \frac{a}{1+a}
    -
    \frac{1}{1+a}
    \left(
        1+\frac{(a^2-1)v^2}{2}
        + \mathcal{O}(v^4)
    \right)\\
    &\approx
    \frac{a-1}{a+1}
    -
    \frac{a-1}{2}v^2.
\end{align*}

This expansion suggests that the leading-order dynamics near $v=0$ can be approximated by an Allen--Cahn-type equation with effective coefficient
\begin{equation*}
    r \approx \frac{a-1}{a+1}.
\end{equation*}

This approximation is further supported by the local slope of the reaction term at the origin. Differentiating $f(v)$ yields
\begin{align*}
    f'(v)
    &=
    v(1-a)v
    \left(
        1+(a^2-1)v^2
    \right)^{-\frac{1}{2}}
    +
    \left(
        1
        -
        \frac{1+\sqrt{1+(a^2-1)v^2}}{1+a}
    \right).
\end{align*}
Evaluating at $v=0$ gives
\begin{equation*}
    f'(0)=\frac{a-1}{a+1}.
\end{equation*}

Similarly, for the Allen--Cahn reaction term
\begin{equation*}
    g(v)=rv(1-v^2),
\end{equation*}
we obtain
\begin{align*}
    g'(v)&=r(1-3v^2),\\
    g'(0)&=r.
\end{align*}
Matching the slopes at the origin therefore leads to the approximation
\begin{equation}
    f(v)\approx rv(1-v^2),
    \qquad
    r=\frac{a-1}{a+1}.
\end{equation}

Consequently, the interaction front can be approximated by the Allen--Cahn traveling-wave profile,
\begin{equation}
    v(z)
    \approx
    \frac{1}{1+e^{\sqrt{2r}\,z}}.
\end{equation}

We next performed numerical simulations of the symmetric system with interaction strength $a=1.5$. The initial condition was chosen such that species $1$ fully occupied the left half of the domain ($x\in(-1,0)$), whereas species $2$ occupied the right half ($x\in(0,1)$). Starting from this sharp interface, a transition region formed near the center of the domain and gradually relaxed toward a stationary profile. The spatial distributions of both species exhibited sigmoid-like shapes and remained symmetric with respect to $x=0$ throughout the simulation.

To quantify convergence, we measured the maximum absolute difference between consecutive simulation steps. The error decreased approximately exponentially over time and continued to decay until the system reached equilibrium, indicating the stability of the transition front. The final distribution of $N_1-N_2$ (obtained when the maximum stepwise error was below $10^{-6}$) was fitted to a logistic function with fixed asymptotic values $K_1=1$, $K_2=-1$, and center $x_0=0$. The logistic approximation showed excellent agreement with the simulated profile. The fitted central slope was $k=0.676$ (after rescaling), which is close to the theoretical prediction
\begin{equation*}
    \hat{k}=\sqrt{2r}=0.632.
\end{equation*}

\begin{figure}[!htbp]
    \centering
    \includegraphics[scale = 1]{symmetric_simu.pdf}
    \caption{
    (A) Spatial dynamics of two symmetric competing species at simulation steps $0$, $11150$, $22300$, $33500$, $44650$, and $55850$. A smooth transition region emerges from the initially sharp boundary and remains stable throughout the simulation. 
    (B) Convergence error, defined as the maximum difference between consecutive simulation steps. The simulation terminates when the error falls below $10^{-6}$. 
    (C) Logistic regression of the final profile of $N_1-N_2$, yielding $k=0.67657$ and $R^2 > 1-10^{-4}$.
    }
    \label{fig:symmetric_simu}
\end{figure}

To further examine the analytical approximation, we analyzed the transformed variables
\begin{equation*}
    u=N_1+N_2,
    \qquad
    v=N_1-N_2.
\end{equation*}
In the theoretical derivation, the total density $u$ was assumed to remain approximately constant and close to $1$. This assumption is supported by the simulations: away from the interface, where only one species dominates, the total density remained close to the carrying capacity. Near the transition region, however, $u$ was slightly lower than $1$, reflecting the reduction in total population density caused by strong interspecific competition. As the simulation evolved, the deviation of $u$ from unity became increasingly localized around the interface.

In contrast, the difference variable $v$ evolved from the initial step-like profile into a smooth sigmoid-shaped distribution, consistent with the predicted logistic traveling-wave form.

\begin{figure}[!htbp]
    \centering
    \includegraphics[scale = 1]{symmetric_simu_sum_diff.pdf}
    \caption{
    Spatial distributions of $v=N_1-N_2$ and $u=N_1+N_2$ during the simulation. Curves correspond to simulation steps $0$, $100$, $200$, $500$, $1000$, $2000$, $5000$, and $10000$. At later times, the profiles converge toward stationary distributions and therefore become visually indistinguishable.
    }
    \label{fig:symmetric_u-v}
\end{figure}

To evaluate the validity of the analytical approximation over different interaction strengths, we further simulated the symmetric system for varying values of
\begin{equation*}
    q=a-1,
\end{equation*}
where $a$ denotes the interspecific competition coefficient. As $q$ increased, the fitted central slope $k$ also increased, indicating that stronger competition produces a narrower transition region and therefore weaker spatial coexistence.

The Allen--Cahn approximation provides a good description only when $a$ is close to $1$. For small $q$, the effective coefficient
\begin{equation*}
    r=\frac{a-1}{a+1}=\frac{q}{q+2}
\end{equation*}
approaches zero, yielding
\begin{equation*}
    \hat{k}=\sqrt{2r}\to0,
\end{equation*}
which is consistent with the broad transition fronts observed numerically. However, when $a\to\infty$, the interface width should vanish and thus $k\to\infty$, whereas the approximation predicts only
\begin{equation*}
    \hat{k}\to\sqrt{2}.
\end{equation*}
This discrepancy indicates that the Allen--Cahn reduction is accurate only in the weak-competition regime and breaks down for strong interspecific interactions.

\begin{figure}[!htbp]
    \centering
    \includegraphics[scale = 0.8]{kq.png}
    \caption{
    Relationship between the simulated central slope $k$ and the theoretical estimate $\hat{k}$ as a function of $q=a-1$. The approximation agrees well with simulations for weak to moderate competition strengths, but systematically underestimates $k$ when $q$ becomes large.
    }
    \label{fig:kq}
\end{figure}

\subsection{Balance Between Dispersal, Reproduction, and Competition}

In more general parameter regimes, the two-species system does not necessarily converge to a stationary transition front. We first investigated the combined effects of diffusion and intrinsic growth by varying the diffusion coefficient $D$ and intrinsic growth rate $r$ over the interval $(0,2)$. Simulations were performed for $100000$ time steps with temporal step size $1\times10^{-6}$. The direction and speed of the transition front were characterized by the position $x^*$ satisfying
\begin{equation*}
    N_1(x^*) = N_2(x^*).
\end{equation*}
A balanced coexistence state corresponds to the condition $x^*=0$, indicating that the interface remains stationary throughout the simulation.

The resulting phase diagram reveals a positive relationship between diffusion rate and intrinsic growth rate along the balance condition. Species with higher dispersal ability therefore require correspondingly larger growth rates in order to maintain a stable spatial coexistence front. This behavior differs fundamentally from the classical locally mixed coexistence condition in the generalized Lotka--Volterra model, where coexistence is determined solely by interaction coefficients.

\begin{figure}[!htbp]
    \centering
    \includegraphics[scale = 1]{scan_param.pdf}
    \caption{
    Transition dynamics and balance conditions under different parameter combinations.
    (A) Phase diagram in the $(r,D)$ plane on linear scales. Blue regions indicate parameter combinations for which species $1$ expands, whereas red regions indicate expansion of species $2$. Lighter colors correspond to slower front velocities. The black curve denotes the balance condition at which the transition front remains stationary.
    (B) The same $(r,D)$ phase diagram represented on logarithmic scales.
    (C) Phase diagram in the $(a,D)$ plane on linear scales.
    (D) The corresponding $(a,D)$ phase diagram on logarithmic scales.
    }
    \label{fig:D-r_balance}
\end{figure}

To further characterize the balance relation between dispersal and reproduction, we analyzed the phase boundary in logarithmic coordinates. The balance condition appears approximately linear in the $\log D$--$\log r$ plane, with slope close to $1/2$. This observation suggests the scaling relation
\begin{equation*}
    D \sim r^2.
\end{equation*}
Therefore, the maintenance of a stationary coexistence front requires diffusion and local population growth to satisfy an approximate quadratic balance relation.

For parameter combinations close to the balance boundary, the two species appear to coexist for long periods of time, although the interface position satisfying $N_1=N_2$ is displaced from the center of the domain. However, long-term simulations reveal that this apparent coexistence is only transient. To illustrate this behavior, we simulated the case $D=1.25$ and $r=0.75$, corresponding to a parameter combination located slightly above the balance line. In this regime, species $1$ eventually invades the entire domain and species $2$ becomes extinct.

The convergence dynamics also differ substantially from those observed in the symmetric case. Initially, the simulation error decreases rapidly as the sharp initial boundary relaxes into a smooth transition front, similar to the balanced symmetric system. After the front is established, however, the error remains approximately constant over an extended period. This plateau corresponds to the slow propagation of the interface across the domain: although the front translates spatially, its shape remains nearly unchanged, resulting in only small differences between consecutive simulation steps.

As the moving front approaches the domain boundary, the population distribution changes rapidly because of the no-flux boundary condition, leading to a temporary increase in the simulation error. Finally, once species $2$ is eliminated, the system converges toward the single-species equilibrium and the error decreases rapidly again.

These dynamics are analogous to those of an asymmetric Allen--Cahn system, in which the two stable states possess unequal potentials and the interface therefore propagates toward the energetically unfavorable state.

\begin{figure}[!htbp]
    \centering
    \includegraphics[scale = 1]{insymme.pdf}
    \caption{
    Simulation of an asymmetric two-species system with $D=1.25$ and $r=0.75$.
    (A) Spatial population dynamics. A transition front first forms and subsequently propagates toward the left boundary, eventually leading to the extinction of species $2$.
    (B) Convergence error during the simulation. The error initially decreases rapidly, then remains near $10^{-6}$ during front propagation, exhibits a transient peak near the boundary, and finally decays after competitive exclusion is completed.
    (C) Temporal evolution of $N_1-N_2$. A logistic-like transition front forms initially and subsequently translates leftward.
    (D) Temporal evolution of the total population density $N_1+N_2$.
    }
    \label{fig:insymme_simu}
\end{figure}

A similar analysis was performed for the balance relation between diffusion rate $D$ and competition strength $a_{1,2}$. In a locally mixed competitive community, species with interspecific competition coefficient smaller than unity are generally excluded by stronger competitors. In contrast, the spatial model reveals that coexistence can still emerge if the weaker competitor disperses sufficiently slowly. Specifically, species with lower competitive ability were able to coexist with a species whose competition coefficient was fixed at $1.5$, provided that their diffusion rate was sufficiently reduced.

Moreover, the balance relation between $D$ and $a_{1,2}$ is also positive: stronger competitive interactions require larger dispersal rates to maintain a stationary coexistence front. In logarithmic coordinates, the phase boundary approximately follows a power-law relation with exponent between $3$ and $4$.

\subsection{A Special Case of Multispecies Competitive Communities}

For multispecies systems, the dynamics become substantially more complex and require a much larger parameter space to be explored numerically. To obtain a qualitative understanding of transition fronts in multispecies communities, we therefore focus on several simplified cases. Although these assumptions are not fully realistic ecologically, they allow the existence and structure of stable transition fronts to be analyzed more clearly.

First, we consider communities with purely competitive interactions. Restricting interactions to competition suppresses oscillatory or chaotic dynamics near equilibrium and constrains the total biomass of the system. Second, we assume that all species are identical in their dispersal and intrinsic growth rates, namely
\begin{equation*}
    D_i = 1,
    \qquad
    r_i = 1.
\end{equation*}
Under these assumptions, we introduce a minimal toy model of multispecies bistability.

To generate bistable community dynamics at the global scale, we partition the species into two competing subcommunities. Species within the same subcommunity interact with strength $a=1$, whereas species belonging to different subcommunities compete more strongly, with interaction strength $b>1$. The resulting interaction matrix takes the form
\begin{equation*}
    \begin{pmatrix}
        1       &\dots&\textcolor{blue}{a}      &\textcolor{red}{b}   &\textcolor{red}{\dots}  &\textcolor{red}{b}  \\
        \dots   &\dots&\dots                    &\textcolor{red}{\dots}  &\textcolor{red}{\dots}  &\textcolor{red}{\dots}  \\
        \textcolor{blue}{a}&\dots&1             &\textcolor{red}{b}   &\textcolor{red}{\dots}  &\textcolor{red}{b}  \\
        \textcolor{red}{b}&\textcolor{red}{\dots} &\textcolor{red}{b}  &1    &\dots&\textcolor{blue}{a}  \\
        \textcolor{red}{\dots}  &\textcolor{red}{\dots}  &\textcolor{red}{\dots}
        &\dots   &\dots &\dots  \\
        \textcolor{red}{b}&\textcolor{red}{\dots} &\textcolor{red}{b}  &\textcolor{blue}{a}   &\dots&1
    \end{pmatrix}.
\end{equation*}

Within each local subcommunity, the generalized Lotka--Volterra dynamics reduce to
\begin{equation}
    \frac{\mathrm{d}N_i}{\mathrm{d}t}
    =
    N_i
    \left(
        1-\sum_{j=1}^{S}N_j
    \right).
\end{equation}
At equilibrium, the total biomass satisfies
\begin{equation*}
    \sum_{j=1}^{S}N_j^*=1,
\end{equation*}
whereas the equilibrium abundance of each individual species depends on the initial condition. Biologically, this degeneracy arises because species within the same subcommunity are dynamically equivalent and therefore behave effectively as different components of the same functional group.

Let $N_{1,i}$ denote species belonging to subcommunity $1$, and $N_{2,j}$ species belonging to subcommunity $2$. The corresponding diffusion--reaction system is
\begin{equation*}
    \begin{aligned}
        \frac{\partial N_{1,i}}{\partial t}
        &=
        \nabla^2 N_{1,i}
        +
        N_{1,i}
        \left(
            1
            -
            \sum_{k}N_{1,k}
            -
            \sum_{\ell}bN_{2,\ell}
        \right),\\
        \frac{\partial N_{2,j}}{\partial t}
        &=
        \nabla^2 N_{2,j}
        +
        N_{2,j}
        \left(
            1
            -
            \sum_{\ell}N_{2,\ell}
            -
            \sum_{k}bN_{1,k}
        \right).
    \end{aligned}
\end{equation*}

Because all species within the same subcommunity experience identical interaction terms, the system can be reduced to the dynamics of the total subcommunity biomasses,
\begin{equation*}
    C_n=\sum_i N_{n,i}.
\end{equation*}
The governing equations become
\begin{equation*}
    \begin{aligned}
        \frac{\partial C_1}{\partial t}
        &=
        \nabla^2 C_1
        +
        C_1
        \left(
            1
            -
            C_1
            -
            bC_2
        \right),\\
        \frac{\partial C_2}{\partial t}
        &=
        \nabla^2 C_2
        +
        C_2
        \left(
            1
            -
            C_2
            -
            bC_1
        \right).
    \end{aligned}
\end{equation*}

This reduced system is mathematically identical to the symmetric two-species model analyzed previously, with the total subcommunity biomasses replacing the species abundances. Consequently, the difference between the two subcommunities is also expected to exhibit an approximately logistic transition profile,
\begin{equation}
    C_1(x)-C_2(x)
    \approx
    \frac{1}{
        1+e^{\sqrt{2\frac{b-1}{b+1}}\,x}
    }.
\end{equation}

Since the abundance of each species remains proportional to the total biomass of its subcommunity, the spatial distributions of individual species are likewise expected to follow approximate logistic forms.

\begin{figure}[!htbp]
    \centering
    \includegraphics[scale = 1]{multi_symme.pdf}
    \caption{
    Simulations of multispecies competitive dynamics.
    (A) Spatial dynamics of a six-species community with intra-community interaction strength $a=1$, inter-community interaction strength $b=1.5$, and identical dispersal and growth rates for all species. A stable transition front forms, and the final species abundances depend on the initial population distribution.
    (B,C) Spatial distributions of $C_1-C_2$ and $C_1+C_2$ during the simulation, showing dynamics analogous to those of the symmetric two-species system.
    (D) Simulation of a sixteen-species community. The two subcommunities remain interactionally symmetric, with the same values of $a$ and $b$ as in panel (A), but species are assigned different dispersal rates:
    $r=1,1.05,1.1,1.15,1.2,1.25,1.3,1.35$ for species $1$--$8$ and $9$--$16$, respectively. The final spatial distribution depends strongly on dispersal ability.
    (E,F) Spatial distributions of $C_1-C_2$ and $C_1+C_2$ corresponding to the simulation shown in panel (D).
    }
    \label{fig:multiple_simu}
\end{figure}

\subsection{Logistic Equilibrium under Strong Inter-Community Competition}

Building upon the simplified toy model described above, we next investigated several perturbed systems that remain close to the symmetric multispecies limit while relaxing some of its assumptions. These simulations were designed to examine the robustness of the logistic transition structure under moderate heterogeneity in interaction strength and dispersal ability.

We first considered a system in which the two local subcommunities remained interactionally symmetric, but the intra-community competition strength was slightly reduced from the neutral case. The system consisted of six species in total, with inter-community competition fixed at $b=1.5$ and intra-community competition set to $a=0.95$. Under these conditions, the transition front remained stationary near the center of the spatial domain because the two subcommunities retained overall symmetry in their ecological properties.

Compared with the case $a=1$, the total equilibrium biomass increased slightly, reflecting the weaker competitive suppression within each subcommunity. Despite this perturbation, logistic regression showed that the equilibrium distributions of all six species were still well approximated by logistic functions. This result suggests that the logistic structure of the coexistence front is robust to moderate deviations from perfectly symmetric competition.

We next investigated the effects of heterogeneous dispersal rates. Relative diffusion coefficients were assigned as
\begin{equation*}
    D_i = 1,\ 1.05,\ 1.1.
\end{equation*}
The transition front again remained stable, indicating that moderate dispersal heterogeneity does not destroy the bistable coexistence structure. However, species with lower dispersal rates occupied narrower spatial regions and exhibited reduced equilibrium abundances, implying that dispersal ability influences both the spatial extent and carrying capacity of individual species within the transition zone.

Logistic regression continued to provide an accurate approximation of the equilibrium population distributions, further supporting the hypothesis that logistic-like transition fronts emerge generically in strongly competitive multispecies diffusion--reaction systems.

\subsection{Measure of $\beta$-Diversity under Logistic Spatial Distributions}

The simulations described above provide a natural framework for measuring biodiversity using continuous spatial distributions rather than discrete community samples. In particular, for bistable coexistence systems with transition fronts, $\beta$-diversity plays an important role in characterizing spatial species turnover across the interface between alternative community states.

To quantify biodiversity in this setting, we extended the classical definitions of $\alpha$-, $\beta$-, and $\gamma$-diversity to continuous population distributions. Specifically, the average local diversity was defined as the spatial average of $\alpha(x)$, whereas $\gamma$-diversity was calculated from the spatially averaged abundance of each species across the entire domain. Using the multispecies simulation described above as an example, we first evaluated the local diversity profile $\alpha(x)$.

\begin{figure}[!htbp]
    \centering
    \includegraphics[scale = 0.8]{alpha.png}
    \caption{Spatial distribution of $\alpha$-diversity across the transition region.}
\end{figure}

The results show that $\alpha$-diversity reaches its maximum near the center of the transition zone. This observation is consistent with the ecological expectation that biodiversity is elevated near the boundary between distinct communities, where species originating from both subcommunities coexist locally. In this sense, the transition front acts as a ``soft'' ecological boundary rather than a sharp interface.

The spatially averaged local diversity was calculated as
\begin{equation*}
    \bar{\alpha}=1.181.
\end{equation*}
The average abundance of each species,
\begin{equation*}
    \bar{N}_i,
\end{equation*}
was then obtained by integrating the population distributions over space and used to compute the global diversity,
\begin{equation*}
    \gamma=1.635.
\end{equation*}
Consequently, the resulting $\beta$-diversity was
\begin{equation*}
    \beta
    =
    \frac{\gamma}{\bar{\alpha}}
    =
    1.384.
\end{equation*}

Based on this framework, we propose a continuous-distribution approach for measuring $\beta$-diversity in spatial ecological systems. Given a set of spatial abundance data, the species distributions can first be approximated by logistic functions, after which biodiversity indices can be evaluated analytically from the corresponding continuous distributions. This approach provides a bridge between spatial reaction--diffusion models and classical biodiversity theory, and may be particularly useful for characterizing gradual community transitions and ecological clines.